% ---- Abstract: 200 words max, unstructured (T-ASE). Budget 0.2 page.
% JAMA abstract was structured (Importance/Objectives/...). Rewritten, not trimmed.
\begin{abstract}
Emergency medical services (EMS) must choose a destination hospital for a suspected stroke patient within minutes, without knowing the stroke subtype, and the choice determines whether time-critical treatments are reached in time. We formulate this dispatch decision as a Markov decision process (MDP) over patient location, time since symptom onset, and subtype uncertainty, with hospitals classified by on-site capability (endovascular thrombectomy, intravenous thrombolysis, or neither), road travel times from a routing engine, inter-facility transfer, and published in-hospital treatment intervals. The MDP is solved online by depth-two forward search, and the resulting policy is distilled into a shallow decision tree that EMS crews can apply without computation. In simulations of \pending{N} patients across 10 replicates in the San Francisco Bay Area, the MDP policy raised the mean probability of a excellent 90-day outcome (modified Rankin Scale 0--1) from \pending{x.xxx} under nearest-hospital routing to \pending{x.xxx}, with the largest gains in areas distant from thrombectomy centres and no losses in well-served areas; a depth-\pending{d} tree recovered \pending{xx}\% of this gain. Results were robust to travel-time error and to door-in-door-out assumptions, and a single routing decision took \pending{xxx}~ms, supporting real-time deployment. A Rhode Island case study shows the policy matches a centralised protocol while distributing load across hospitals.
\end{abstract}

% ---- Note to Practitioners: 100-300 words, REQUIRED, immediately after the abstract.
% Audience: practising engineers and EMS medical directors outside stroke research.
% Must cover: practical problem, what the method gives them, near-term use, limitations.
\section*{Note to Practitioners}
\new{This paper was motivated by a routing problem that EMS agencies face every day: an ambulance crew with a suspected stroke patient must pick a hospital in minutes, while the information that determines the right choice (the stroke subtype) is unavailable until the patient reaches a scanner. Most regions resolve this with a fixed rule, such as nearest hospital or nearest comprehensive centre. We show how to compute the destination that maximises the expected probability of a good outcome, accounting for travel times, each hospital's treatment capability, the possibility of an onward transfer, and the time already elapsed since onset. The computed policy is then approximated by a decision tree with a handful of branches on quantities a dispatcher already has (time since onset and drive times to the nearest centres of each type), so it can be printed on a card or embedded in a dispatch system without any optimisation at run time. In two U.S. regions the tree recovers almost all of the benefit of the full optimisation. The hospital lists and travel-time engine are inputs, so the method transfers to a new region by changing data, not code, and can be re-run when hospital designations change. Limitations: the model assumes hospitals have capacity for the patients routed to them, uses one set of in-hospital treatment intervals per hospital type, and has been evaluated in simulation rather than against observed patient outcomes; a prospective evaluation with an EMS agency is the natural next step.}
